% =========================================================
%  PLANTILLA DE MATRIZ DE RIESGOS
%  Actividad A03 - Administración de Riesgos
%  Seguridad Aplicada a Sistemas de Información
%  Compilar con: pdfLaTeX (Overleaf o TeX Live/MiKTeX)
% =========================================================
\documentclass[11pt,a4paper]{article}

% ----------------- Paquetes -----------------
\usepackage[utf8]{inputenc}
\usepackage[spanish]{babel}
\usepackage{geometry}
\geometry{left=2.2cm,right=2.2cm,top=2.5cm,bottom=2.5cm}
\usepackage[table]{xcolor}
\usepackage{tcolorbox}
\tcbuselibrary{skins,breakable}
\usepackage{enumitem}
\usepackage{tabularx}
\usepackage{booktabs}
\usepackage{titlesec}
\usepackage{fancyhdr}
\usepackage{hyperref}
\usepackage{graphicx}
\usepackage{multirow}
\usepackage{array}

% ----------------- Colores institucionales -----------------
\definecolor{azulmat}{RGB}{0,84,140}
\definecolor{griscaja}{RGB}{245,245,245}
\definecolor{azulsuave}{RGB}{232,240,250}
\definecolor{amarillosuave}{RGB}{255,255,224}
% Colores de la matriz de calor
\definecolor{nivelbajo}{RGB}{144,238,144}
\definecolor{nivelmedio}{RGB}{255,235,120}
\definecolor{nivelalto}{RGB}{255,170,90}
\definecolor{nivelcritico}{RGB}{255,90,90}

\hypersetup{
  colorlinks=true,
  linkcolor=azulmat,
  urlcolor=azulmat,
  citecolor=azulmat
}

% ----------------- Cajas de contenido -----------------
\newtcolorbox{instruccion}[1][Instrucciones]{%
  enhanced, breakable,
  colback=amarillosuave, colframe=orange!75!black,
  fonttitle=\bfseries, title=#1
}

\newtcolorbox{informacion}[1][Información]{%
  enhanced, breakable,
  colback=azulsuave, colframe=azulmat,
  fonttitle=\bfseries, title=#1
}

% Caja vacía para completar (texto libre)
\newtcolorbox{campolibre}[1][2cm]{%
  enhanced, breakable,
  colback=white, colframe=griscaja!80!black,
  height=#1, valign=top, boxrule=0.5pt
}

% ----------------- Encabezado y pie -----------------
\pagestyle{fancy}
\fancyhf{}
\fancyhead[L]{\small Matriz de Riesgos --- A03}
\fancyhead[R]{\small Seguridad Aplicada a Sistemas de Información}
\fancyfoot[C]{\thepage}
\renewcommand{\headrulewidth}{0.4pt}

% ----------------- Formato de títulos -----------------
\titleformat{\section}{\large\bfseries\color{azulmat}}{\thesection.}{0.6em}{}
\titleformat{\subsection}{\normalsize\bfseries\color{azulmat!80!black}}{\thesubsection.}{0.6em}{}

% ----------------- Portada -----------------
\title{%
  \rule{\textwidth}{1pt}\\[0.4cm]
  \textbf{\Large Matriz de Riesgos}\\[0.2cm]
  \large Actividad A03 --- Administración de Riesgos\\[0.1cm]
  \normalsize Análisis y Gestión de Riesgos en Sistemas de Información\\[0.2cm]
  \rule{\textwidth}{1pt}\\[0.4cm]
  \normalsize Seguridad Aplicada a Sistemas de Información\\
  Licenciatura en Sistemas de Información --- Facultad de Informática y Diseño\\
  Docente: Ing.\ Rodrigo Atilio Elgueta\\[0.2cm]
  \small Ciclo Académico Vigente
}
\author{}
\date{}

\begin{document}
\maketitle
\thispagestyle{fancy}
\tableofcontents
\newpage

% =========================================================
\section{Datos Generales}\label{sec:datos}

\begin{center}
\renewcommand{\arraystretch}{1.6}
\begin{tabularx}{\textwidth}{@{} p{5cm} X @{}}
\toprule
\textbf{Nombre de la empresa / organización} & Clínica privada (caso desarrollado para TP SimpleRisk) \\
\midrule
\textbf{Rubro / Industria} & Salud --- atención médica privada \\
\midrule
\textbf{Tamaño (empleados / facturación)} & 120 empleados, aprox. 800 pacientes atendidos por día \\
\midrule
\textbf{Alumno(s) / Grupo} & Nicolás Estrella \\
\midrule
\textbf{Legajo(s)} & 31.066.692 \\
\midrule
\textbf{Fecha de elaboración} & 10/09/2026 \\
\midrule
\textbf{Responsable del análisis} & Nicolás Estrella \\
\bottomrule
\end{tabularx}
\end{center}

% =========================================================
\section{Objetivo y Alcance del Análisis}

\begin{campolibre}[2.5cm]
Identificar, evaluar y proponer el tratamiento de los principales riesgos de seguridad de la
información que afectan a una clínica privada, con foco en la confidencialidad, integridad y
disponibilidad de las historias clínicas electrónicas (HCE) y los datos de pacientes, obras
sociales y facturación.
\end{campolibre}

\textbf{Alcance:} (sistemas, procesos, sedes o áreas incluidas en el análisis)

\begin{campolibre}[2cm]
El análisis abarca el servidor de historias clínicas electrónicas, el proceso de facturación
tercerizado con obras sociales, los accesos del personal administrativo y de enfermería al
sistema de HCE, y la infraestructura eléctrica y de red del sitio principal de atención.
No incluye otras sedes ni sistemas administrativos ajenos a la gestión clínica.
\end{campolibre}

% =========================================================
\section{Inventario y Clasificación de Activos}

Asigná un código único a cada activo (ej.: A01, A02\ldots). La clasificación de la información puede ser: \textit{Pública, Interna, Confidencial, Restringida}.

\begin{center}
\small
\renewcommand{\arraystretch}{1.6}
\begin{tabularx}{\textwidth}{@{} >{\centering\arraybackslash}p{1cm} X >{\raggedright\arraybackslash}p{2.2cm} >{\raggedright\arraybackslash}p{2.3cm} >{\raggedright\arraybackslash}p{2.3cm} >{\centering\arraybackslash}p{1.8cm} @{}}
\toprule
\textbf{ID} & \textbf{Descripción del Activo} & \textbf{Tipo} & \textbf{Responsable} & \textbf{Clasificación} & \textbf{Criticidad} \\
\midrule
A01 & Servidor de Historias Clínicas Electrónicas (HCE) & Software / Hardware & Nicolás Estrella & Restringida & Alta \\
A02 & Base de datos de pacientes (datos de salud) & Información & Nicolás Estrella & Restringida & Alta \\
A03 & Datos de obras sociales y facturación & Información & Nicolás Estrella & Confidencial & Alta \\
A04 & Sistema de facturación tercerizado & Software & Proveedor externo & Confidencial & Media \\
A05 & Credenciales de acceso del personal & Información & Nicolás Estrella & Restringida & Alta \\
A06 & Infraestructura eléctrica (UPS / tablero) & Hardware & Nicolás Estrella & Interna & Media \\
A07 & Red interna de la clínica & Red & Nicolás Estrella & Interna & Media \\
A08 & Personal administrativo y de enfermería & Humano & Nicolás Estrella & --- & Media \\
A09 & Copias de seguridad (backups) & Información & Nicolás Estrella & Restringida & Alta \\
A10 & Imagen institucional / reputación de la clínica & Imagen & Dirección & --- & Media \\
\bottomrule
\end{tabularx}
\end{center}

\textit{Tipos de activo: Información / Software / Hardware / Red / Humano / Imagen. Criticidad: Baja / Media / Alta.}

% =========================================================
\section{Identificación de Amenazas y Vulnerabilidades}

Asigná un código único a cada amenaza (ej.: T01, T02\ldots). El tipo puede ser: \textit{Natural, Accidental, Intencional}.

\begin{center}
\small
\renewcommand{\arraystretch}{1.6}
\begin{tabularx}{\textwidth}{@{} >{\centering\arraybackslash}p{1cm} >{\centering\arraybackslash}p{1.8cm} X >{\raggedright\arraybackslash}p{2.2cm} X @{}}
\toprule
\textbf{ID} & \textbf{Activo (ID)} & \textbf{Descripción de la Amenaza} & \textbf{Tipo} & \textbf{Vulnerabilidad Asociada} \\
\midrule
T01 & A01 & Ataque de ransomware que cifra el servidor de HCE & Intencional & Sin EDR ni backups aislados \\
T02 & A02 & Acceso indebido de personal sin necesidad de conocer & Intencional & Falta de control de acceso basado en roles \\
T03 & A09 & Pérdida de datos por falla de hardware o error humano & Accidental & Backups inexistentes / nunca probados \\
T04 & A03, A04 & Filtración de datos por brecha en el proveedor externo & Intencional & Sin cláusulas contractuales de seguridad ni auditorías \\
T05 & A01, A06 & Corte de energía prolongado sin respaldo & Accidental & UPS insuficiente, sin generador \\
T06 & A05, A08 & Phishing dirigido a personal administrativo & Intencional & Falta de capacitación y de MFA \\
T07 & A02, A03 & Incumplimiento normativo (Ley 25.326) & Accidental & Sin procedimientos formales de resguardo de datos \\
\bottomrule
\end{tabularx}
\end{center}

% =========================================================
\section{Escalas de Valoración}\label{sec:escalas}

\subsection{Escala de Probabilidad}

\begin{center}
\renewcommand{\arraystretch}{1.4}
\begin{tabularx}{\textwidth}{@{} c p{2.8cm} X @{}}
\toprule
\textbf{Valor} & \textbf{Nivel} & \textbf{Descripción} \\
\midrule
1 & Raro & El evento solo ocurriría en circunstancias excepcionales. \\
\midrule
2 & Improbable & Podría ocurrir, pero no se espera que suceda. \\
\midrule
3 & Posible & Existe una posibilidad real de que ocurra. \\
\midrule
4 & Probable & Es muy probable que ocurra en algún momento. \\
\midrule
5 & Casi seguro & Se espera que ocurra frecuentemente / ha ocurrido recientemente. \\
\bottomrule
\end{tabularx}
\end{center}

\subsection{Escala de Impacto}

\begin{center}
\renewcommand{\arraystretch}{1.4}
\begin{tabularx}{\textwidth}{@{} c p{2.8cm} X @{}}
\toprule
\textbf{Valor} & \textbf{Nivel} & \textbf{Descripción} \\
\midrule
1 & Insignificante & Impacto mínimo, sin consecuencias operativas ni económicas relevantes. \\
\midrule
2 & Menor & Alteración leve de la operación, costo bajo. \\
\midrule
3 & Moderado & Impacto operativo y económico apreciable, recuperable. \\
\midrule
4 & Mayor & Impacto significativo, pérdida de operación o datos, costo alto. \\
\midrule
5 & Catastrófico & Paralización total, pérdida crítica de datos, daño reputacional o legal severo. \\
\bottomrule
\end{tabularx}
\end{center}

% =========================================================
\section{Matriz de Calor (Probabilidad × Impacto)}

La matriz 5×5 relaciona la \textbf{probabilidad} (eje vertical) con el \textbf{impacto} (eje horizontal). Cada celda muestra el valor de riesgo resultante (\(P \times I\)) y su color indica el nivel.

\begin{center}
\renewcommand{\arraystretch}{1.9}
\setlength{\tabcolsep}{8pt}
\begin{tabular}{@{} c c | c c c c c @{}}
\multicolumn{2}{c}{} & \multicolumn{5}{c}{\textbf{IMPACTO} $\longrightarrow$} \\
\multicolumn{2}{c}{} & \textbf{1} & \textbf{2} & \textbf{3} & \textbf{4} & \textbf{5} \\
\cline{3-7}
\multirow{5}{*}{\rotatebox[origin=c]{90}{\textbf{PROBABILIDAD} $\longrightarrow$}}
 & \textbf{5} & \cellcolor{nivelmedio}5 & \cellcolor{nivelalto}10 & \cellcolor{nivelalto}15 & \cellcolor{nivelcritico}20 & \cellcolor{nivelcritico}25 \\
 & \textbf{4} & \cellcolor{nivelbajo}4 & \cellcolor{nivelmedio}8 & \cellcolor{nivelalto}12 & \cellcolor{nivelcritico}16 & \cellcolor{nivelcritico}20 \\
 & \textbf{3} & \cellcolor{nivelbajo}3 & \cellcolor{nivelmedio}6 & \cellcolor{nivelmedio}9 & \cellcolor{nivelalto}12 & \cellcolor{nivelalto}15 \\
 & \textbf{2} & \cellcolor{nivelbajo}2 & \cellcolor{nivelbajo}4 & \cellcolor{nivelmedio}6 & \cellcolor{nivelmedio}8 & \cellcolor{nivelalto}10 \\
 & \textbf{1} & \cellcolor{nivelbajo}1 & \cellcolor{nivelbajo}2 & \cellcolor{nivelbajo}3 & \cellcolor{nivelbajo}4 & \cellcolor{nivelmedio}5 \\
\cline{3-7}
\end{tabular}
\end{center}

\subsection{Niveles de Riesgo}\label{sec:niveles}

\begin{center}
\renewcommand{\arraystretch}{1.4}
\begin{tabularx}{\textwidth}{@{} p{3cm} p{2.5cm} p{3.5cm} X @{}}
\toprule
\textbf{Nivel} & \textbf{Rango (P×I)} & \textbf{Color} & \textbf{Acción requerida} \\
\midrule
\cellcolor{nivelbajo} \textbf{Bajo} & 1 -- 4 & Verde & Monitorear. Puede aceptarse. \\
\cellcolor{nivelmedio} \textbf{Medio} & 5 -- 9 & Amarillo & Plan de acción a mediano plazo. \\
\cellcolor{nivelalto} \textbf{Alto} & 10 -- 15 & Naranja & Tratamiento prioritario a corto plazo. \\
\cellcolor{nivelcritico} \textbf{Crítico} & 16 -- 25 & Rojo & Acción inmediata. Escalar a dirección. \\
\bottomrule
\end{tabularx}
\end{center}

% =========================================================
\section{Evaluación de Riesgos}\label{sec:evaluacion}

Asigná un código único a cada riesgo (ej.: R01, R02\ldots). Calculá el valor como \(Probabilidad \times Impacto\) y clasificá el nivel según la tabla anterior.

\begin{center}
\small
\renewcommand{\arraystretch}{1.6}
\begin{tabularx}{\textwidth}{@{} >{\centering\arraybackslash}p{0.8cm} >{\centering\arraybackslash}p{1.2cm} >{\centering\arraybackslash}p{1.6cm} >{\centering\arraybackslash}p{1.1cm} >{\centering\arraybackslash}p{1.3cm} >{\centering\arraybackslash}p{1.1cm} >{\centering\arraybackslash}p{1.4cm} X @{}}
\toprule
\textbf{ID} & \textbf{Activo} & \textbf{Amenaza} & \textbf{Prob.} & \textbf{Impacto} & \textbf{Valor} & \textbf{Nivel} & \textbf{Justificación breve} \\
\midrule
R01 & A01 & T01 & 4 & 5 & 20 & \cellcolor{nivelcritico}Crítico & Sector salud muy atacado por ransomware; sin EDR ni backups aislados. \\
R02 & A02 & T02 & 4 & 4 & 16 & \cellcolor{nivelcritico}Crítico & Sin control de acceso por rol; visibilidad total de HCE. \\
R03 & A09 & T03 & 3 & 5 & 15 & \cellcolor{nivelalto}Alto & Sin política formal de backups ni pruebas de restauración. \\
R04 & A05, A08 & T06 & 4 & 4 & 16 & \cellcolor{nivelcritico}Crítico & Sin capacitación ni MFA; puerta de entrada a otros ataques. \\
R05 & A03, A04 & T04 & 3 & 4 & 12 & \cellcolor{nivelalto}Alto & Proveedor externo sin cláusulas de protección de datos. \\
R06 & A01, A06 & T05 & 4 & 3 & 12 & \cellcolor{nivelalto}Alto & Cortes de energía frecuentes, sin respaldo suficiente. \\
R07 & A02, A03 & T07 & 3 & 4 & 12 & \cellcolor{nivelalto}Alto & Sin procedimientos documentados conforme a Ley 25.326. \\
\bottomrule
\end{tabularx}
\end{center}

% =========================================================
\section{Tratamiento de Riesgos y Riesgo Residual}\label{sec:tratamiento}

Para cada riesgo, definí la estrategia de tratamiento:
\begin{itemize}[leftmargin=*]
  \item \textbf{Mitigar:} implementar salvaguardas para reducir probabilidad o impacto.
  \item \textbf{Transferir:} pasar el riesgo a un tercero (seguros, outsourcing).
  \item \textbf{Aceptar:} asumir el riesgo conscientemente.
  \item \textbf{Evitar:} discontinuar la actividad que lo genera.
\end{itemize}

\begin{center}
\small
\renewcommand{\arraystretch}{1.6}
\begin{tabularx}{\textwidth}{@{} >{\centering\arraybackslash}p{1.3cm} >{\raggedright\arraybackslash}p{1.7cm} X >{\raggedright\arraybackslash}p{1.7cm} >{\centering\arraybackslash}p{1cm} >{\centering\arraybackslash}p{1cm} >{\centering\arraybackslash}p{1cm} >{\centering\arraybackslash}p{1.3cm} @{}}
\toprule
\textbf{Riesgo} & \textbf{Estrategia} & \textbf{Salvaguardas propuestas} & \textbf{Tipo salv.} & \textbf{Prob. resid.} & \textbf{Imp. resid.} & \textbf{Val. resid.} & \textbf{Nivel resid.} \\
\midrule
R01 & Mitigar & Backups automatizados aislados/inmutables + EDR + segmentación de red & Técnica & 2 & 4 & 8 & \cellcolor{nivelmedio}Medio \\
R02 & Mitigar & Control de acceso basado en roles (RBAC) + auditoría de accesos & Técnica & 2 & 4 & 8 & \cellcolor{nivelmedio}Medio \\
R03 & Mitigar & Backups automatizados diarios con pruebas de restauración periódicas & Técnica & 2 & 5 & 10 & \cellcolor{nivelalto}Alto \\
R04 & Mitigar & Capacitación periódica + simulacros de phishing + MFA obligatorio & Administrativa & 2 & 3 & 6 & \cellcolor{nivelmedio}Medio \\
R05 & Transferir & Cláusulas contractuales de protección de datos + auditorías al proveedor & Administrativa & 2 & 3 & 6 & \cellcolor{nivelmedio}Medio \\
R06 & Mitigar & UPS de mayor capacidad + evaluación de generador de respaldo & Física & 2 & 3 & 6 & \cellcolor{nivelmedio}Medio \\
R07 & Mitigar & Documentar procedimientos y designar responsable de tratamiento de datos & Administrativa & 2 & 3 & 6 & \cellcolor{nivelmedio}Medio \\
\bottomrule
\end{tabularx}
\end{center}

\textit{Tipo de salvaguarda: Técnica / Física / Administrativa. El nivel residual se clasifica con la misma tabla de la Sección~\ref{sec:niveles}.}

% =========================================================
\section{Resumen de Resultados}

Distribución de los riesgos identificados por nivel:

\begin{center}
\renewcommand{\arraystretch}{1.4}
\begin{tabularx}{0.85\textwidth}{@{} p{4cm} p{3.5cm} X @{}}
\toprule
\textbf{Nivel} & \textbf{Cantidad de riesgos} & \textbf{\% del total} \\
\midrule
\cellcolor{nivelcritico} Crítico & 3 & 43\% \\
\cellcolor{nivelalto} Alto & 4 & 57\% \\
\cellcolor{nivelmedio} Medio & 0 & 0\% \\
\cellcolor{nivelbajo} Bajo & 0 & 0\% \\
\midrule
\textbf{Total} & 7 & 100\% \\
\bottomrule
\end{tabularx}
\end{center}

\section{Conclusiones y Recomendaciones}

\begin{campolibre}[4cm]
El análisis identificó 7 riesgos relevantes, de los cuales 3 resultaron de nivel Crítico
(ransomware sobre el servidor de HCE, acceso indebido a historias clínicas y phishing al
personal administrativo) y 4 de nivel Alto. La causa común detrás de la mayoría de estos
riesgos es la ausencia de controles técnicos básicos: backups aislados, control de acceso por
roles y autenticación multifactor. Se recomienda priorizar la implementación de backups
inmutables y MFA como medidas de bajo costo y alto impacto, dado que reducen sensiblemente
la probabilidad de los riesgos más críticos. Tras aplicar las salvaguardas propuestas, ningún
riesgo residual queda por encima del nivel Alto, lo que valida la estrategia de tratamiento
definida.
\end{campolibre}

\section{Declaración de Buenas Prácticas}

\begin{informacion}[Compromiso profesional]
Como estudiante de Seguridad Aplicada a Sistemas de Información, declaro que el presente análisis de riesgos fue elaborado aplicando las buenas prácticas de la disciplina, con criterio ético, fundamentos técnicos y respetando el marco normativo vigente.
\end{informacion}

\vspace{1.2cm}

\begin{center}
\begin{tabularx}{\textwidth}{@{} X X @{}}
\rule{6cm}{0.4pt} & \rule{6cm}{0.4pt} \\
Firma del alumno/a & Fecha \\
\end{tabularx}
\end{center}

% =========================================================
\section{Referencias Bibliográficas}

\begin{itemize}[leftmargin=*]
  \item Chicano Tejada, E. (2023). \textit{Auditoría de seguridad informática. IFCT0109} (2.ª ed.). IC Editorial --- eLibro.
  \item ISO/IEC (2005). \textit{Information security management 27001}.
  \item ISO/IEC (2018). \textit{ISO/IEC 27005: Information security risk management}.
  \item NIST (2012). \textit{Guide for Conducting Risk Assessments (SP 800-30 Rev. 1)}.
  \item OWASP Risk Rating Methodology: \url{https://owasp.org/www-community/OWASP_Risk_Rating_Methodology}
\end{itemize}

\end{document}
